% Research candidate, not integrated or numbered.
% Authored 2026-11-20 after checkpoint21. No priority claim.
\subsection*{Candidate: relative homotopy normalization of the RP3 cylinder test}

Let $B=\mathbb{RP}^3$, $N=B\times I$, and let $P_{\rm CAT}(B)$
denote concordances fixing $B\times\{0\}$, with the top free.
Put
\[
 \mathcal H=\{F\in\operatorname{Homeo}(N\operatorname{ rel}\partial N):
       F\simeq\operatorname{id}_N\operatorname{ rel}\partial N\}.
\]
Here the displayed homotopy is a homotopy of maps, not an isotopy.

\textbf{Candidate assertion.}\quad
Every component of $P_{\rm TOP}(B)$ contains a representative in
$\mathcal H$. Moreover, the equality
\[
 \operatorname{im}\bigl(\pi_0P_{\rm PL}(B)\longrightarrow
                  \pi_0P_{\rm TOP}(B)\bigr)=\ker c_B
 \tag{RH.1}
\]
is equivalent to the following concrete relative assertion:
every $F\in\mathcal H$ with $\operatorname{cs}(F)=0$ is TOP-isotopic,
relative to the entire $\partial N$, to a PL homeomorphism of $N$.
The PL endpoint then automatically belongs to $\mathcal H$.
This is a normalization of the remaining problem, not a proof of
either side of (RH.1).

\textbf{Original calculation used.}\quad
Jahren--Kwasik, \emph{Free involutions on $S^1\times S^n$},
Math. Ann. 351 (2011), 281--303, DOI 10.1007/s00208-010-0599-y,
Proposition 3.5 and Remark 3.6(ii), printed pp.~294--296, give
\[
 \pi_1\operatorname{hAut}(B)
   =\operatorname{im}\pi_1\operatorname{PSO}(4)\oplus\mathbb Z/2.
 \tag{RH.2}
\]
The first summand is naturally the one represented by loops of
projective orthogonal transformations. The extra generator gives a
homotopy self-equivalence of $S^1\times B$ with nonzero TOP normal
invariant. All the other summands of
$\pi_0\operatorname{hAut}(S^1\times B)$ have representatives which
are homeomorphisms. The proof on printed p.~296 describes the extra
generator by the pinch followed by the element of $\pi_4(B)$.
The complete published PDF was opened at
\texttt{https://link.springer.com/content/pdf/10.1007/s00208-010-0599-y.pdf};
retrieval references were \texttt{turn249438676view0} and
\texttt{turn127561356view0}. No structure-set calculation is being
substituted for an isotopy classification.

\begin{proof}
\textbf{First fix the top.}\quad
For $F\in P_{\rm TOP}(B)$, its top map $f$ is homotopic to the
identity, since both inclusions $B\hookrightarrow B\times I$ are
homotopy equivalences and $F$ fixes the bottom. It is orientation
preserving. The projective Smale theorem and the ordinary
three-dimensional DIFF--TOP comparison, already checked in A.162,
show that $f$ is isotopic to the identity. Choose a path $f_u$ from
the identity to $f$, and set
\[
 T_s(x,t)=(f_{st}(x),t),\qquad 0\le s\le1.
\]
Then $T_s$ is a TOP isotopy in $P_{\rm TOP}(B)$, with $T_0$ the
identity and $T_1|_{B\times\{1\}}=f$. Replace $F$ by $T_1^{-1}F$.
This preserves its concordance component and fixes both ends.

\textbf{Remove the relative homotopy class.}\quad
Write the resulting map as $F(x,t)=(u(x,t),q(x,t))$ and put
$a_t(x)=u(x,t)$. Every $a_t$ is a homotopy equivalence: it is the
composition of the level inclusion into $N$, the homeomorphism $F$,
and the homotopy equivalence $N\to B$. Thus $a$ is a loop based at
the identity in $\operatorname{hAut}(B)$. The formula
\[
 (x,t)\longmapsto\bigl(u(x,t),(1-v)q(x,t)+vt\bigr),
       \qquad 0\le v\le1,
 \tag{RH.3}
\]
is a homotopy relative to both ends from $F$ to the level map
$A(x,t)=(a_t(x),t)$.

Glue the two ends of $N$ by the identity. Since $F$ fixes them,
it descends to an actual homeomorphism $\widehat F$ of $S^1\times B$.
Formula (RH.3) descends to a homotopy from $\widehat F$ to
$\widehat A$. Consequently the TOP normal invariant of
$\widehat A$ is zero. In (RH.2), its extra $\mathbb Z/2$ coordinate
must vanish. Indeed, composition with a homeomorphism carries a
nonzero normal invariant to a nonzero normal invariant, whereas
all remaining generators are represented by homeomorphisms.
Hence $[a]$ lies in the projective orthogonal summand.

Use the ordinary automorphism comparison in dimension three from
A.134 to choose a based, jointly PL loop $b_t$ of PL automorphisms
of $B$ representing $[a]$ in $\pi_1\operatorname{hAut}(B)$.
Its trace $T_b(x,t)=(b_t(x),t)$ is a PL homeomorphism fixing both
ends. It is null in $P_{\rm TOP}(B)$, by the explicit isotopy
$(x,t)\mapsto(b_{st}(x),t)$. The map
\[
             G=T_b^{-1}F
\]
therefore represents the same concordance component as $F$.
Postcompose (RH.3) with $T_b^{-1}$. The resulting level loop is
$b_t^{-1}a_t$, which is based null-homotopic in
$\operatorname{hAut}(B)$ by the choice of $b$. This gives
$G\simeq\operatorname{id}_N$ relative to both ends. Thus $G\in\mathcal H$.

\textbf{Free-top isotopy can be made relative for a PL endpoint.}\quad
Suppose $F\in\operatorname{Homeo}(N\operatorname{ rel}\partial N)$
belongs to the PL image on $\pi_0P$. Choose an isotopy in
$P_{\rm TOP}(B)$ from $F$ to a PL concordance $L$.
The PL top map of $L$ is orientation preserving, so the
$\pi_0$ part of A.134 provides a jointly PL isotopy of that map
to the identity. Composing $L$ with its PL trace gives a PL map
$L_0$ fixing both ends, in the same free-top TOP component.
Let $H_s$ be a free-top TOP isotopy from $F$ to $L_0$ and write
$h_s=H_s|_{B\times\{1\}}$; this is a based loop in
$\operatorname{Homeo}(B)$. Put
\[
 Q_s(x,t)=(h_{st}^{-1}(x),t).
\]
Then $Q_sH_s$ is an isotopy relative to both ends from $F$ to
$T_h^{-1}L_0$. The $\pi_1$ surjectivity in A.134 gives a based
jointly PL loop $k$ and a based homotopy of loops from $h$ to $k$.
Taking inverse traces of this homotopy and composing with $L_0$
is an isotopy relative to both ends from $T_h^{-1}L_0$ to
$T_k^{-1}L_0$. The latter is PL. This argument asserts no PL
regularity of the intermediate TOP isotopies.

Now (RH.1) implies the asserted relative statement by this
endpoint correction and isotopy invariance of $c_B$. Conversely,
given a zero-$c_B$ component, the first two steps supply
$G\in\mathcal H$ in that component. For this both-end-fixed
representative, A.143 identifies $c_B([G])$ with its relative
formal Casson--Sullivan class. The proposed relative statement
then puts the component in the PL image. The reverse containment
in (RH.1) is already A.143. This proves the equivalence.
\end{proof}

\textbf{Scope.}\quad
The normal invariant is used only to remove the extra relative
homotopy-equivalence coordinate. Its vanishing is not asserted
to classify either isotopies or PL structures. The final
relative TOP isotopy-to-PL statement remains unproved. No new
numbered theorem, original-row closure, or qualifying movement
is claimed by this candidate alone.
